\section{Reproduction and Additional Details}
\label{sec:appendix:exp}

\subsection{Hardware Specification}
\label{sec:appendix:hwspec}

All experiments run on a single GPU:
\begin{itemize}[leftmargin=1.5em]
  \item \textbf{NVIDIA RTX PRO 6000 Blackwell}: ${\sim}$188 SMs (GB202),
    96\,GB GDDR7 (${\sim}$1.8\,TB/s memory bandwidth),
    5th-generation Tensor Cores (HMMA BF16/FP8), Blackwell architecture (sm\_120).
    Note: the \texttt{warp\_cycles\_per\_instruction} counter is absent from the
    Blackwell counter set; \texttt{eligible\_\allowbreak cycles\_pct} serves as the
    latency-bound indicator instead.
\end{itemize}

Software requirements: NVIDIA GPU (Ampere minimum), \nsys $\ge$ 2024.6,
\ncu $\ge$ 2025.4.1, PyTorch 2.11+, and CUDA 12.8.
Reproduction scripts, \texttt{profile.json} pairs, and per-example counter comparisons are provided in the accompanying repository here: 
\url{https://github.com/yottalabsai/Profiler}

\subsection{Hardware Counter Reference}
\label{sec:appendix:counters}

\begin{table}[h]
\centering
\caption{The 20 hardware counters collected by \sys, their bottleneck axis, and
what they measure.  ``Agg.'' denotes the aggregation method
  applied when multiple kernels are attributed to one operator (see \secref{sec:aggregation}).
  $\dagger$~\texttt{warp\_cycles\_per\_instr} is absent from the Blackwell counter set (sm\_120);
  \texttt{eligible\_\allowbreak cycles\_pct} (counter~14) serves as the latency-bound indicator on the
  RTX PRO 6000 used for all experiments reported here.}
\label{tab:counters}
\scriptsize
\setlength{\tabcolsep}{4pt}
\begin{tabular}{clllc}
\toprule
\textbf{\#} & \textbf{Profile Field} & \textbf{Bottleneck Axis} & \textbf{Measures} & \textbf{Agg.}\\
\midrule
1  & \texttt{gpu\_time\_duration\_ns}      & Latency          & GPU wall time                & SUM \\
2  & \texttt{dram\_bytes\_read}            & Memory BW        & DRAM read bytes              & SUM \\
3  & \texttt{dram\_bytes\_written}         & Memory BW        & DRAM write bytes             & SUM \\
4  & \texttt{memory\_throughput\_pct}      & Memory BW        & Memory subsystem \% peak     & DW-mean \\
5  & \texttt{dram\_throughput\_pct}        & Memory BW        & DRAM \% of peak BW           & DW-mean \\
6  & \texttt{mem\_busy\_pct}              & L1 pipeline      & L1 pipeline \% peak active   & DW-mean \\
7  & \texttt{l1\_hit\_rate}               & Cache            & L1 cache hit rate (\%)       & DW-mean \\
8  & \texttt{l2\_hit\_rate}               & Cache            & L2 cache hit rate (\%)       & DW-mean \\
9  & \texttt{sm\_throughput\_pct}         & Compute          & SM \% of peak throughput     & DW-mean \\
10 & \texttt{tensor\_core\_active\_pct}   & Tensor Cores     & TC pipeline active (\%)      & DW-mean \\
11 & \texttt{sm\_active\_cycles}          & Compute          & Total SM active cycles       & SUM \\
12 & \texttt{achieved\_occupancy}         & Occupancy        & Warps active \% of peak      & DW-mean \\
13$\dagger$ & \texttt{warp\_cycles\_per\_instr}    & Latency          & Cycles per issued instruction & DW-mean \\
14 & \texttt{eligible\_\allowbreak cycles\_pct}       & Latency          & Eligible warp issue rate (\%) & DW-mean \\
15 & \texttt{executed\_instructions}      & Throughput       & Total executed instructions  & SUM \\
16 & \texttt{ipc\_active}                 & Throughput       & IPC when SM active           & DW-mean \\
17 & \texttt{avg\_threads\_per\_warp}     & Divergence       & Avg threads active per warp  & DW-mean \\
18 & \texttt{registers\_per\_thread}      & Reg.\ pressure   & Registers allocated / thread & MAX \\
19 & \texttt{local\_memory\_spills}       & Reg.\ spills     & Spill events to L1 local mem & SUM \\
20 & \texttt{dynamic\_smem\_per\_block}   & SMEM pressure    & Dynamic shared mem / block   & MAX \\
\bottomrule
\end{tabular}
\end{table}

\subsection{Prototype Automation via Claude Code Plugin}
\label{sec:appendix:plugin}

As a prototype implementation, \sys includes a Claude Code plugin that automates the
closed-loop protocol via a single \texttt{/optimize workload.py} command.  The plugin orchestrates five
specialized LLM agents---\textbf{capture-agent}, \textbf{optimization-strategist},
\textbf{backend-engineer}, \textbf{validation-agent}, and \textbf{/report}---covering all
pipeline stages from baseline capture through cross-profile comparison. 


\subsection{Annotated \texttt{profile.json} Schema}
\label{sec:appendix:schema}

\texttt{profile.json} is the central output of \sys: a JSON document that maps each
dispatched PyTorch operator to its constituent CUDA kernels, per-kernel \nsys durations,
and aggregated hardware metrics.  \tabref{tab:schema_fields} describes every top-level field; the listing below
shows a representative single-operator excerpt from the GPT-2 baseline
(\texttt{aten::mm}, 512$\times$768 $\times$ 768$\times$768 weight projection, NVTX-attributed).
The \texttt{aggregated} block condenses per-kernel counter values into a single
operator-level record using the duration-weighted mean (\equref{eq:dwmean}); these are the
values consumed by downstream analysis tools and the Claude Code plugin.
The \texttt{duration\_ns} and \texttt{aggregated.total\_duration\_ns} fields are
\nsys-derived timing fields; \ncu metrics, including
\texttt{gpu\_\_time\_duration.sum}, are stored only inside \texttt{metrics.raw}.

\begin{lstlisting}[basicstyle=\scriptsize\ttfamily, breaklines=true, frame=single,
                   xleftmargin=2pt, xrightmargin=2pt]
{
  "schema_version": "1.0",
  "capture_metadata": {
    "model_name": "gpt2",
    "device_name": "NVIDIA RTX PRO 6000 Blackwell",
    "compile_mode": "inductor",
    "capture_timestamp_utc": "2026-06-01T00:53:25Z"
  },
  "operators": [
    {
      "operator_id": "aten::mm_26",
      "operator_name": "aten::mm",
      "call_index": 26,
      "is_fused": false,
      "fused_with": [],
      "kernels": [
        {
          "kernel_id": "k_00880",
          "kernel_name": "cutlass_80_simt_sgemm_128x32_8x5_nn_align1",
          "attribution_method": "nvtx",
          "confidence": "medium",
          "duration_ns": 32736,
          "stream_id": 7,
          "grid_dim": [128, 1, 5],
          "block_dim": [128, 1, 1],
          "metrics": { "...": "20 raw ncu counter fields (architecture-specific names)" }
        }
      ],
      "aggregated": {
        "total_duration_ns": 32736,
        "kernel_count": 1,
        "tensor_core_active_pct": 0.0,
        "achieved_occupancy": 18.26,
        "sm_throughput_pct": 37.6,
        "dram_throughput_pct": 7.16,
        "l2_hit_rate": 89.55,
        "registers_per_thread": 80.0,
        "eligible_cycles_pct": 40.6,
        "local_memory_spills": 0,
        "warp_cycles_per_instruction": null
      }
    }
  ],
  "unattributed_kernels": []
}
\end{lstlisting}

Key fields: \texttt{attribution\_method} records the tier that assigned operator identity
(\texttt{"nvtx"}, \texttt{"inductor\_fusion"}, or \texttt{"torch\_profiler\_correlation"});
\texttt{confidence} mirrors the tier level (\texttt{"high"} or \texttt{"medium"});
\texttt{is\_fused} and \texttt{fused\_with} list all \texttt{aten::} ops fused into the
same Triton kernel.  The \texttt{null} value for \texttt{warp\_cycles\_per\_instruction}
reflects the absent Blackwell counter.
The \texttt{metrics.raw} sub-object stores all architecture-specific ncu column names
exactly as emitted by \texttt{ncu --import --csv}; the \texttt{aggregated} block outlines an operator-level metric summary of the attributed kernels.

\begin{table}[h]
\centering
\caption{Top-level \texttt{profile.json} field reference.}
\label{tab:schema_fields}
\scriptsize
\setlength{\tabcolsep}{4pt}
\begin{tabular}{lp{3.5cm}p{7.5cm}}
\toprule
\textbf{Field} & \textbf{Type} & \textbf{Meaning} \\
\midrule
\texttt{schema\_version}       & string        & Schema revision for forward-compatibility checks \\
\texttt{capture\_metadata}     & object        & Device name, torch/CUDA versions, nsys/ncu report paths \\
\texttt{operators[]}           & array         & One entry per attributed operator dispatch \\
\quad\texttt{operator\_id}     & string        & Unique key: \texttt{aten::op\_N} \\
\quad\texttt{operator\_name}   & string        & Canonical \texttt{aten::} name (with size annotations) \\
\quad\texttt{call\_index}      & int           & Dispatch ordinal within the measured window \\
\quad\texttt{is\_fused}        & bool          & True if kernel spans $>$1 aten op (Inductor fusion) \\
\quad\texttt{fused\_with}      & string[]      & Co-fused \texttt{aten::} op names \\
\quad\texttt{kernels[]}        & array         & Per-kernel records with raw ncu metrics \\
\quad\texttt{aggregated}       & object        & Duration-weighted operator-level metric summary \\
\texttt{unattributed\_kernels[]} & array       & Kernels matching no attribution path \\
\bottomrule
\end{tabular}
\end{table}

\subsection{Supplementary Figures}
\label{sec:appendix:candidates}

\begin{figure}[H]
\centering
\begin{tikzpicture}[font=\small, >=Stealth,
  stage/.style={draw, rounded corners=4pt, fill=white,
                minimum width=3.6cm, minimum height=0.75cm,
                align=center, inner sep=5pt, font=\small},
  hook/.style={draw, rounded corners=3pt, align=left,
               inner sep=5pt, font=\scriptsize},
  arr/.style={-Stealth,thick},
  harr/.style={-Stealth, thick, dashed},
]

  % Pipeline stages (vertical, left column)
  \node[stage, fill=gray!10] (in)     at (0,0)    {\texttt{workload.py}};
  \node[stage, fill=blue!8]  (dynamo) at (0,-2.2)  {Dynamo FX\\graph capture};
  \node[stage, fill=blue!8]  (aot)    at (0,-4.4)  {AOTAutograd\\decomposition};
  \node[stage, fill=blue!8]  (ind)    at (0,-6.6)  {Inductor\\code generation};
  \node[stage, fill=gray!10] (out)    at (0,-8.8)  {Triton / CUDA kernels};

  \draw[arr] (in)     -- (dynamo);
  \draw[arr] (dynamo) -- (aot);
  \draw[arr] (aot)    -- (ind);
  \draw[arr] (ind)    -- (out);

  % Level 1 hook (between Dynamo and AOT)
  \node[hook, fill=red!8, draw=red!50] (l1) at (5.5,-2.2)
    {\textbf{Level 1} (pre-decomposition)\\
     \textbullet\ QKV fusion (\textit{SDPA case study})\\
     \textbullet\ Attention canonicalization};
  \draw[harr, red!60] (l1.west) -- (dynamo.east);

  % Level 2 hook (between AOT and Inductor)
  \node[hook, fill=orange!8, draw=orange!50] (l2) at (5.5,-4.4)
    {\textbf{Level 2} (Aten IR)\\
     \textbullet\ BF16 dtype promotion (\textit{all case studies})\\
     \textbullet\ Weight folding, layout annotation};
  \draw[harr, orange!60] (l2.west) -- (aot.east);

  % Level 3 hook (at Inductor)
  \node[hook, fill=green!8, draw=green!50!black] (l3) at (5.5,-6.6)
    {\textbf{Level 3} (Inductor config)\\
     \textbullet\ Weight freezing (\textit{all case studies})\\
     \textbullet\ Autotuning mode selection};
  \draw[harr, green!60!black] (l3.west) -- (ind.east);

  % Compose label
  \node[font=\scriptsize, gray, align=center] at (0,-9.6)
    {Passes at each level compose\\without interfering across levels};

\end{tikzpicture}
\caption{\textbf{Three-level FX pass hook points.}
  The \texttt{torch.compile()} pipeline lowers a model through three IR stages.
  Level~1 passes must run before AOTAutograd because structural patterns such as
  a shared \texttt{LayerNorm} output feeding Q, K, and V projections---the
  precondition for QKV fusion---are decomposed into per-consumer copies by
  AOTAutograd and become invisible afterward.  Level~2 (Aten IR) applies
  dtype and layout transforms on the fully decomposed graph.  Level~3
  Inductor flags control non-graph behaviors such as weight freezing and
  autotuning that cannot be expressed as graph rewrites.}
\label{fig:cand:fxpasslevel}
\end{figure}
